% SPDX-License-Identifier: MIT
% fastq_stream -- Application Note draft for Bioinformatics (Oxford Academic).
%
% SUBMISSION NOTE
% ---------------
% Bioinformatics Application Notes use OUP's `bioinfo.cls`, obtainable from the
% journal's LaTeX template pack. That class is not vendored here (licence), so
% this draft targets `article` with a two-column layout matching the published
% dimensions, and compiles with any stock TeX installation. To submit:
%   1. replace the preamble down to \begin{document} with the bioinfo.cls one,
%   2. \documentclass{bioinfo}, \copyrightyear, \pubyear, \access, \appnotes,
%   3. keep the section structure below unchanged -- it already matches.
%
% LIMIT: 4 printed pages, ~2,600 words including references.
%
% STATUS: the fastp/Trimmomatic comparison in Table 2 is NOT YET MEASURED.
% Those cells read TODO on purpose. Do not submit with estimated numbers in
% them; run scripts/run_benchmarks.py on GIAB HG002 and paste the generated
% markdown table. Every number that IS present in this draft was measured on
% the host described in Section 3.1.

\documentclass[twocolumn,10pt]{article}

\usepackage[margin=2cm]{geometry}
% Deliberately no `times`: it remaps \ttfamily to Courier (pcr), whose metrics
% are absent from minimal TeX installations, and this draft is dense with
% \texttt. bioinfo.cls sets the submission fonts itself, so the dependency buys
% nothing and costs portability. Default CM fonts compile everywhere.
\usepackage{booktabs}
% natbib for author-year \citep, matching bioinfo.cls and the \bibitem[...]
% author-year forms used in the bibliography below.
\usepackage[round]{natbib}
\usepackage{graphicx}
\usepackage{url}
\usepackage{amsmath}
\usepackage[colorlinks=true,allcolors=blue]{hyperref}
\usepackage{listings}

\lstset{basicstyle=\ttfamily\scriptsize,breaklines=true,frame=none,
        columns=fullflexible,xleftmargin=0pt}

\setlength{\parskip}{0pt}
\title{\vspace{-2em}\textbf{fastq\_stream: lock-free in-memory FASTQ streaming\\
       for disk-free NGS preprocessing}}
\author{\small [Author list]\\[-0.3em]
        \small \textit{[Affiliation]}}
\date{}

\begin{document}
\maketitle

% ============================================================ ABSTRACT
\noindent\textbf{Abstract}

\noindent\textbf{Summary:} Next-generation sequencing preprocessing is
routinely performed by shell pipelines that decompress a gzipped FASTQ, trim it,
recompress it, write it to disk, and read it back into an aligner. The
arithmetic of quality control is negligible against the cost of that round trip.
We present \texttt{fastq\_stream}, a C++20 engine that performs single-pass
quality control, adapter removal and quality trimming entirely in memory and
streams clean reads directly into an aligner through a named pipe, eliminating
the intermediate compressed FASTQ altogether. Stages communicate through
genuinely single-producer/single-consumer lock-free ring buffers with
cache-line-isolated indices, achieving a median inter-stage hand-off latency of
14\,ns. Quality kernels are vectorised for AVX2/AVX-512 and NEON and are
asserted bit-identical to scalar references at every read length. Peak resident
memory is approximately 10\,MB and is independent of input size.

\noindent\textbf{A structural finding:} we show that parallel gzip
decompression is impossible for the single-member gzip framing that sequencers
emit, and that re-framing FASTQ as BGZF at generation time yields a
2.5$\times$ decompression speedup for every downstream consumer, at no cost.

\noindent\textbf{Availability and implementation:} C++20 source, CMake build,
and container definitions at \url{https://github.com/[org]/fastq_stream}, MIT
licensed.

\noindent\textbf{Contact:} \texttt{[email]}

% ============================================================ INTRO
\section{Introduction}

A sequencing core facility producing terabytes of gzipped FASTQ per day
typically preprocesses it with a pipeline of the form

\begin{lstlisting}
zcat in.fq.gz | trimmomatic ... | gzip > tmp.fq.gz
bwa-mem2 mem ref.fa tmp.fq.gz > out.sam
\end{lstlisting}

Three costs dominate, and none is the quality-control arithmetic itself.
First, \emph{serial decompression}: \texttt{zlib} inflate sustains roughly
200--400\,MB/s on one core. Second, \emph{intermediate materialisation}: the
trimmed FASTQ is DEFLATE-compressed (several times more expensive than
inflation), written to disk, read back, and inflated again. Third,
\emph{interpreter and pipe overhead} from the Bash or Python glue.

Existing tools attack parts of this. \texttt{fastp}
\citep{chen2018fastp} is a fast C++ implementation with
multithreading and integrated QC reporting; \texttt{Trimmomatic}
\citep{bolger2014trimmomatic} is the widely used Java reference for
window-based trimming; \texttt{seqtk} provides minimal C primitives. All of
them, in standard use, still write a compressed intermediate.
\texttt{libdeflate} \citep{biggers2016libdeflate} offers substantially
faster inflate and deflate than \texttt{zlib}, and \texttt{pigz} parallelises
compression but not decompression of an existing stream.

\texttt{fastq\_stream} removes the intermediate entirely, reduces decompression
cost where the input framing permits, and replaces pipe glue with in-process
lock-free hand-offs.

% ============================================================ METHODS
\section{Implementation}

\subsection{Pipeline topology}

The engine is a five-rank pipeline: a reader, an optional pool of $D$
inflaters, a record assembler, a pool of $W$ quality-control workers, and a
writer. Chunks of 256\,KiB flow between ranks; only buffer handles move, and
payload bytes are never copied between stages.

\subsection{Genuinely SPSC queues}

A single-producer/single-consumer ring is correct only with exactly one
producer thread and one consumer thread. A shared queue drained by a worker
pool is multi-producer/multi-consumer and requires different, slower
algorithms. Rather than weaken the guarantee, the topology is arranged so that
no ring is ever touched by more than two threads: the reader owns one lane per
inflater, each inflater owns one output lane, the assembler owns one lane per
worker, and each worker owns one output lane.

Each ring has power-of-two capacity with one slot sacrificed so that
$\mathit{head}=\mathit{tail}$ unambiguously denotes emptiness. The head and
tail indices are separated by \texttt{std::hardware\_destructive\_interference\_size}
to eliminate false sharing, and each side caches the opposite index in
thread-private storage so that a steady-state operation never reads the other
core's cache line. Measured hand-off latency between two threads is 14\,ns at
70\,M items/s. The implementation is clean under ThreadSanitizer.

\subsection{Ordering without a reorder buffer}

Chunk $k$ is dispatched to lane $k \bmod D$, and each lane preserves order
internally. A collector reading lanes round-robin in the same order therefore
recovers the original stream order with no sequence-number sorting and no
reorder buffer. The single requirement this imposes is that a worker emit
exactly one output buffer per input chunk, which holds because trimming only
shrinks records and identifier lines are copied verbatim.

Because FASTQ permits \texttt{@} as a quality character, record boundaries
cannot be located by scanning for \texttt{@}. The assembler instead counts
newlines from a known boundary and cuts only on multiples of four lines,
splicing any straddling record into 64\,KiB of headroom reserved in front of
the next chunk.

\subsection{Vectorised quality kernels}

Per-cycle quality sums, $Q\!\ge\!20$ and $Q\!\ge\!30$ fractions, and adapter
Hamming extension are vectorised with compile-time dispatch to AVX-512BW, AVX2
or NEON. Per-cycle sums accumulate into 32-bit lanes folded into 64-bit totals
once per chunk; a 256\,KiB chunk holds approximately 650 reads, so overflow
would require $5.8\times10^{7}$ reads at a single cycle. Every vectorised path
is asserted bit-identical to a scalar reference at all lengths from 0 to 300,
covering the sub-vector tails where such kernels typically fail.

\subsection{Memory}

Buffers are 64-byte aligned and recycled through a pool. The pool free list is
mutex-guarded, which is not a contradiction of the lock-free claim: it is
touched once per 256\,KiB of stream, against hundreds of millions of lock-free
ring operations. Peak resident memory is approximately 10\,MB and is
structurally bounded by
$(\mathrm{lanes}\times\mathrm{slots}+\mathrm{in\ flight})\times 256\,\mathrm{KiB}$,
independent of input size.

% ============================================================ RESULTS
\section{Results}

\subsection{Experimental setup}

Component and end-to-end measurements below were obtained on an Apple
M-series host, 12 cores, NEON backend, on 150\,bp simulated reads with a
NovaSeq-like quality decay profile. This is a hybrid-core laptop: consecutive
runs of an unchanged kernel varied by up to a factor of two through core
placement and clock drift, so every figure reported is the best of $N$ trials
($N=5$ for kernels, $N=3$ end-to-end). These numbers establish ranking and
scaling behaviour; they are not a substitute for the cluster measurements in
Section~3.4.

\subsection{Input framing determines decompression parallelism}

\texttt{libdeflate} exposes no streaming interface: it inflates a complete gzip
member into a caller-sized buffer. This is not an API oversight but the source
of its speed, and it has a structural consequence. A \emph{single-member} gzip
stream --- what \texttt{bcl2fastq}, DRAGEN and \texttt{gzip} produce --- has one
member spanning the whole file, and no interior DEFLATE block boundary can be
located without inflating everything before it. Parallel decompression of such
a stream is therefore impossible, for any tool. BGZF, by contrast, is a
sequence of independent members each declaring its compressed size and
expanding to at most 64\,KiB, so members can be located by header arithmetic
and inflated concurrently.

\texttt{fastq\_stream} implements both paths and reports which was taken
(Table~\ref{tab:framing}). Output is byte-identical between them.

\begin{table}[h]
\centering
\caption{Effect of input framing. 200k $\times$ 150\,bp reads; identical output.}
\label{tab:framing}
\small
\begin{tabular}{lrr}
\toprule
Framing & Throughput & Inflater threads \\
\midrule
BGZF (\texttt{bgzip}) & 914\,MB/s & 3 \\
Plain gzip (single member) & 372\,MB/s & 0 \\
\bottomrule
\end{tabular}
\end{table}

The practical implication is disproportionate to its difficulty: running
\texttt{bgzip} rather than \texttt{gzip} once, at the point of FASTQ
generation, gives a 2.5$\times$ decompression speedup to every downstream
consumer of that file, permanently. Any benchmark comparing tools must state
the framing used; reporting a BGZF number against a competitor's plain-gzip
number would be invalid.

\subsection{Scaling and component throughput}

End-to-end throughput was 645, 1863, 3063 and 4304\,MB/s at 1, 3, 5 and 9
workers respectively --- $6.7\times$ from $9\times$ the workers, saturating as
memory bandwidth and the serial assembler stage become binding. Peak buffer
pool occupancy remained under 83\,MB throughout.

Component throughput (best of 5): \texttt{sum\_u8} 55\,GB/s, \texttt{count\_ge}
49\,GB/s, per-cycle accumulation 43\,GB/s, sliding-window trimming 89\,GB/s.
Adapter seed-and-extend runs at 0.95\,GB/s and is thus the pipeline's compute
bottleneck, roughly fifty times more costly than any quality kernel. Replacing
a variable-length \texttt{memcmp} in the seed scan with a single unaligned
32-bit comparison improved it by $3.4\times$.

\subsection{Comparison against existing tools}

\begin{table}[h]
\centering
\caption{GIAB HG002 NovaSeq 2$\times$150\,bp. \textbf{NOT YET MEASURED} ---
see submission note at the head of this file.}
\label{tab:baselines}
\small
\begin{tabular}{lrrr}
\toprule
Tool & Wall (s) & Peak RSS (MB) & Disk written (GB) \\
\midrule
\texttt{fastq\_stream} (FIFO) & TODO & TODO & TODO \\
\texttt{fastq\_stream} (file) & TODO & TODO & TODO \\
\texttt{fastp}                & TODO & TODO & TODO \\
\texttt{Trimmomatic}          & TODO & TODO & TODO \\
\bottomrule
\end{tabular}
\end{table}

These cells are deliberately unpopulated. \texttt{fastp} and
\texttt{Trimmomatic} were not available on the development host, and the
project's benchmark harness (\texttt{scripts/run\_benchmarks.py}) has been
written but not yet executed against them. The harness isolates each
measurement in a separate process, because \texttt{ru\_maxrss} over
\texttt{RUSAGE\_CHILDREN} is a monotonic high-water mark that otherwise leaks
the peak of the heaviest tool into every subsequent measurement; it reports
medians over replicates; and it measures \texttt{fastq\_stream} on both
framings so the comparison against a plain-gzip baseline is like for like.

\subsection{Adapter specificity}

Seed-and-extend matching trades sensitivity for speed, and the cost is
measurable. On 200k random reads with 40{,}000 planted TruSeq adapters, the
conventional minimum overlap of 4 produced 3{,}528 spurious trims (2.2\% of
reads), falling to 248 (0.15\%) at an overlap of 8. We retain 4 as the default
for consistency with the cutadapt and fastp families but recommend 8 for
production use, which removes 93\% of spurious calls while missing only adapter
remnants shorter than 8\,bp. Separately, a read whose adapter copy carries a
sequencing error within the seed window is missed entirely --- the same
compromise \texttt{fastp} makes.

% ============================================================ CONCLUSION
\section{Conclusion}

\texttt{fastq\_stream} demonstrates that the dominant cost in FASTQ
preprocessing is the compressed intermediate file, not the quality-control
computation, and that removing it is straightforward given a strictly
single-producer/single-consumer streaming architecture. The most transferable
result is not the implementation but Section~3.2: the ubiquitous practice of
emitting single-member gzip FASTQ forecloses parallel decompression for every
downstream tool, and the remedy costs one flag at generation time.

Current limitations are single-end input only, a 64\,KiB maximum record size
(adequate for short-read but not long-read data), uncompressed output by
design, and the absence of BAM input.

% ============================================================ REFS
\begin{thebibliography}{9}
\small
\bibitem[Bolger \textit{et~al.}(2014)]{bolger2014trimmomatic}
Bolger,A.M., Lohse,M. and Usadel,B. (2014) Trimmomatic: a flexible trimmer for
Illumina sequence data. \textit{Bioinformatics}, \textbf{30}, 2114--2120.

\bibitem[Chen \textit{et~al.}(2018)]{chen2018fastp}
Chen,S., Zhou,Y., Chen,Y. and Gu,J. (2018) fastp: an ultra-fast all-in-one
FASTQ preprocessor. \textit{Bioinformatics}, \textbf{34}, i884--i890.

\bibitem[Eggers(2016)]{biggers2016libdeflate}
Biggers,E. (2016) libdeflate: heavily optimized library for DEFLATE/zlib/gzip
compression and decompression. \url{https://github.com/ebiggers/libdeflate}.

\bibitem[Li(2013)]{li2013bwa}
Li,H. (2013) Aligning sequence reads, clone sequences and assembly contigs with
BWA-MEM. \textit{arXiv}:1303.3997.

\bibitem[Zook \textit{et~al.}(2016)]{zook2016giab}
Zook,J.M. \textit{et~al.} (2016) Extensive sequencing of seven human genomes to
characterize benchmark reference materials. \textit{Sci. Data}, \textbf{3},
160025.
\end{thebibliography}

\end{document}
